\section{Coprime scale carriers and fixed prime separation}
\label{sec:r13-radicals}

A low-degree scale and fixed prime separation can be combined without
adding a variable. The resulting complete integer equations are
\begin{center}
\begin{tabular}{lrrr}\toprule
Construction & Variables & Exact degree & Binary operations\\\midrule
Normalized five-root norm &13&1072&192\\
Ordinary four-root norm &14&378&145\\
Ordinary three-root norm &15&172&123\\
Centered three-root norm &15&170&134\\
Centered two-root norm &16&98&138\\\bottomrule
\end{tabular}
\end{center}
All variables range over unrestricted integers. These are explicit
cost--degree tradeoffs, not global minima. In particular, the first row
improves the incoming $13/1088/195$ point; the second substantially lowers
both degree and evaluation cost from the inspected $14/652/163$ point.
The $16$-variable row lowers degree but increases cost relative to
$16/108/132$.

The source is the literal signed CRT parent
\texttt{degree88\_17.json}, with its inherited ordinary native compiler
and K215 calibration. We restore each of its five residuals below.
This is not a new Lean/Rocq certificate, and it proves neither consistency
nor a historical or global optimality claim.

\subsection{Complete common prefix}
The base thirteen coordinates are
\[
 j,C,z,h,r,v,w,T,c,n,t,p_2,p_3.
\]
The fixed six numerals are the coherent source-34 tuple
$\beta,K_0,M_C,M_F$, with $M_F$ already shifted, and
$A_{\rm cone}=2^{1113}$, $B_{\rm cone}=2^{666104}$.
In particular $\beta=2^{1064921}-1$. All named registers below are
polynomial abbreviations, not extra quantified variables:
\begin{align*}
J&=6j+1,&q&=\beta J+1,&W&=12w+4,\\
X&=qW,&Y&=qv,&E&=XY,\\
a&=E+Y,&H&=4a+3,&\Delta&=a^2+H,\\
F&=(K_0+W)C-z(q-1),&&&\\
R&=(q(q-F)-C)(q^2-1)+(M_C+qM_F)J,&&&\\
k&=R+1+hE,&\eta&=c-Yk,&\zeta&=k-\eta,\\
u&=EYk,&A_\tau&=u(u+k)+1,&A_m&=1+\Delta c^2,\\
b_0&=X+ac,&I&=(2r+1)c^2,&M&=2q^2,\\
U&=MH,&e&=v-q^2.&&
\end{align*}
The determinant and guard radicand are
\begin{align*}
d&=(A_{\rm cone}F)(B_{\rm cone}C)
 -(q-A_{\rm cone}F-B_{\rm cone}C)^2,\\
A_g&=d\eta\zeta-(t^2-d-\eta\zeta+1)^2
     -1-p_2^2-p_3(p_3+1).
\end{align*}
For the ordinary variants, add the coordinate $f$ and set
\[
 Q=I^2,\qquad S=Q-\Delta(f^2-1),\qquad
 V=fcT-f^2R-c.
\]
For the thirteen-variable normalized variant instead use
\[
 A_f=1+\Delta I^2,\qquad Q=\Delta^2I^2,\qquad
 V=A_f(cT-R)-c.
\]
In both cases put $\delta=Q-1$, $C_y=QV^2-1$. The rational expression
$A_y=C_y/\delta$ is proof notation only, never a circuit division.
Where the coordinate $\gamma$ is retained, write
$D=(b_0+H\gamma)^2-A_m$. Where $p_1$ is retained, write
$G=A_g-p_1^2$.

\subsection{Five literal norm specifications}
Let
\[
 \mathcal F_m(x;a_1,\ldots,a_m)
 =\prod_{\varepsilon\in\{-1,1\}^m}
   \left(x+\sum_i\varepsilon_i\sqrt{a_i}\right).
\]
Sign invariance proves that this is an integer polynomial. The following
expressions are cleared symbolically; their JSON circuits use only binary
addition, subtraction and multiplication.

For thirteen variables take $Z=U(n+V)+Mb_0+He$ and
\[
 N_{13}=\delta^{16}\mathcal F_5\left(
 Z;81U^2A_f,9U^2A_\tau,M^2A_m,729U^2A_g,
 \frac{U^2C_y}{\delta}\right).
\]
The equation is $N_{13}=0$.

For fourteen variables, add $f$, retain the same $Z$, and use
\[
 N_{14}=\delta^8\mathcal F_4\left(
 Z;9U^2A_\tau,81U^2A_g,M^2A_m,
 \frac{U^2C_y}{\delta}\right),\qquad
 P_{14}=(1+N_{14})(1+S)-1.
\]
For fifteen variables, add $f,\gamma$, and use
\[
 N_{15}=\delta^4\mathcal F_3\left(
 M(n+V)+e;9M^2A_\tau,81M^2A_g,
 \frac{M^2C_y}{\delta}\right),
 \qquad P_{15}=(1+N_{15})(1+S)(1+D)-1.
\]

The centered variants use $t_0=2u+k$, $\ell=k^2-4$.
For fifteen variables add $f,p_1$ and set
\begin{align*}
 Z_c&=U(n+3t_0+2V)+2Mb_0+He,\\
 N_{15c}&=\delta^4\mathcal F_3\left(
 Z_c;36U^2A_\tau,4M^2A_m,\frac{4U^2C_y}{\delta}\right),\\
 P_{15c}&=(1+N_{15c})(1+S)(1+4G)-1.
\end{align*}
For sixteen variables add $f,\gamma,p_1$, put $x=Mn+e$, and define
\begin{align*}
 j_0&=\delta\{x^2+M(6t_0+4V)x+9M^2\ell\}
       -4M^2(V^2-1),\\
 h_0&=12\delta t_0VM^2,\\
 N_{16}&=j_0(j_0+2h_0)
  +144\delta M^4\{4A_\tau(1-V^2)+\delta\ell V^2\},\\
 P_{16}&=(1+N_{16})(1+S)(1+D)(1+4G)-1.
\end{align*}
The exact identity, valid also at $\delta=0$, is
\[
 N_{16}=\delta^2\mathcal F_2\left(
 M(n+3t_0+2V)+e;36M^2A_\tau,
 \frac{4M^2C_y}{\delta}\right)
\]
after symbolic clearing. An independent integer-coefficient check compares
all $29$ nonzero terms in the centered identity.

\subsection{Unit products recover the rows before denominator use}
For integers, if residuals $R_i$ never equal $-2$, then
\[
 (1+N)\prod_i(1+R_i)=1
 \quad\Longleftrightarrow\quad N=R_1=\cdots=R_s=0.
\]
Indeed every factor of the integer unit $1$ is $\pm1$. Each $R_i$ factor
is $+1$, forcing the one unrestricted norm factor to be $+1$ too.
There is no missing assumption on the norm modulo four.

Here $a$ is even and $\Delta\equiv3\pmod4$. If $f$ is odd then
$S\equiv I^2\pmod4$; if $f$ is even then $S\equiv I^2-1\pmod4$.
Thus $S$ is never $2$ modulo four. The same is true of $D$: for even $c$
its residues are $0,3$, and for odd $c$ they are $0,1$.
Finally $4G\equiv0\pmod4$. All the stated finalizers therefore recover
their individual residuals, in particular $S=0$, before interpreting a
rational norm.

Since $a\in6\mathbb Z$, $\Delta=(a+1)(a+3)\ge3$.
In ordinary variants $S=0$ excludes $Q=1$, as it would give
$\Delta(f^2-1)=1$. In the normalized variant $Q=\Delta^2I^2$ excludes
$Q=1$ directly. Hence $\delta\ne0$ at every zero in question.

\subsection{Prime separation before guard recovery}
The fixed residues give
\[
 q\equiv X\equiv2\pmod3,\quad a\equiv H\equiv\Delta\equiv0\pmod3,
 \quad A_\tau\equiv A_m\equiv1\pmod3.
\]
To check $A_\tau$, write $A_\tau=1+b(b+1)k^2$ with
$b=XY^2\equiv0$ or $2$ modulo three.
Also $H$ is nonzero and odd. Under $S=0$, the square $Q=I^2$ is divisible
by three, hence by nine. The normalized $Q$ is already divisible by nine.
Thus $A_y$ is a three-adic unit of residue one.

The weighted roots used by the ordinary fourteen-variable norm have
squared orders
\[
 -2v_3(H),\quad0,\quad2,\quad4+v_3(A_g)
\]
for $d_0/H,y,3\tau,9p_1$ respectively. The normalized construction
inserts $9f$ of squared order four and uses $27p_1$ of squared order
at least six. A zero guard root is harmless.

Here is the precise separation argument. In any subset of at least two
nonzero roots, the smallest squared order is unique and even. Divide
out its power of three and multiply the putative zero sum over all signs
of the remaining roots. The product is a polynomial in the rational
squares and is a nonzero unit modulo three, a contradiction.
Subtracting a Galois conjugate of a rational signed root sum would give
exactly such a forbidden signed subsum. Every root is therefore rational.
This argument allows negative radicands; their rationality subsequently
rules that case out. No positivity guard is assumed beforehand.

Roots of integer squares are integral. For the auxiliary root, $Q=0$
gives $y=\pm1$. For $Q>1$, rational $y$ in
$QV^2-(Q-1)y^2=1$ is integral by adjacent-Pell descent. In lowest terms
$y=A/m$ gives $m^2\mid Q-1$. The transformation
\[
 V'=(2Q-1)V-2(Q-1)y,\qquad
 y'=(2Q-1)y-2QV
\]
preserves the norm and the denominator, and after taking positive signs
satisfies $0<V'<V$ when $V>1$. The endpoint $|V|=1$ forces $m=1$.
Thus the case $c=0$, and consequently $Q=0$, has not been silently excluded.

\subsection{Coprime carriers and exact source restoration}
Because $Y=qv$, we have $H\equiv3\pmod q$. Since $q\equiv2\pmod3$
and $H$ is odd,
\[
 \gcd(H,2q^2)=1.
\]
After root recovery the ordinary fourteen-variable factor is
\[
 U(n+V)+Mb_0+He=Md_0+U(3\tau+y+9p_1).
\]
Reduction modulo $M$ and $H$ gives respectively
\[
 v=q^2(2v_{\rm old}+1),\qquad
 \gamma=\frac{d_0-b_0}{H}\in\mathbb Z.
\]
The normalized case adds $9f$ and replaces $9p_1$ by $27p_1$;
the same argument applies. The centered main carrier instead has
$2M(d_0-b_0)$ modulo $H$, still invertible because $H$ is odd.
The variants retaining $\gamma$ need only the first congruence.
In all cases $Y=(2v_{\rm old}+1)q^3$ exactly, including its odd quotient.
The removed $v^2$ guard multiplier is unnecessary: $v\ne0$ follows
now from the recovered congruence, not from an assumed inequality.

The integral guard root restores $G=0$, which implies $d>0$ and
$\eta\zeta>0$ by the standard joint-positivity identity. For completeness,
if both factors were negative their magnitudes $b,e>0$ would require
$(t^2+b+e+1)^2<be$, which is impossible. Thus $c=0$ is impossible too,
since it would give $\eta\zeta=-Y(Y+1)k^2\le0$.

For the ordinary variants set old $J_3=2j$, old scale $v_{\rm old}$
as above, and restore the two centered source coordinates by
\[
 \tau_{\rm old}=2\tau-(2u+k),\qquad y_{\rm gap}=y-V.
\]
The first and auxiliary parent residuals become exactly
\begin{align*}
 \tau_{\rm old}^2+2\tau_{\rm old}(2u+k)+k^2-4
   &=4(\tau^2-A_\tau),\\
 V^2+(1-Q)y_{\rm gap}(2V+y_{\rm gap})-1
   &=QV^2-(Q-1)y^2-1.
\end{align*}
The other three rows are the main norm, $S$ and $G$.
Every remaining parent abbreviation agrees literally. This restores the
five-row signed CRT source directly, without inventing a new rank lemma.
For normalized thirteen variables instead set
\[
 r_{\rm old}=\frac{\Delta(2r+1)-1}{2},\qquad
 T_{\rm old}=fT.
\]
The first quotient is integral because $\Delta$ is odd. Then
$I_{\rm old}=\Delta I$, the parent's $V$ agrees exactly, and its strong
row follows from $f^2=A_f$.

Conversely, the inherited canonical refutation witnesses have odd
$J\equiv1\pmod3$, so $J=6j+1$. Put
$v=q^2(2v_{\rm old}+1)$ and retain their guard witnesses.
For example the ordinary fourteen-variable construction chooses
\[
 n=\gamma+3\tau+y+9p_1-V-e/M.
\]
The centered fifteen-variable construction chooses
\[
 n=6\tau+2y+2\gamma-3t_0-2V-e/M.
\]
The other formulas omit the coordinates not fused by their norms.
All values are integers and make one conjugate factor vanish.
The normalized thirteen-variable case uses the same canonical odd strong
normalization and fresh auxiliary reconstruction proved in the incoming
prime-separation package. Its ordinary calibration boundary is unchanged.

\subsection{Universal degree bounds and actual coefficient certificates}
Fix the six ports. The exact propagated prefix bounds are
\[
\begin{array}{c|rrrrrrrrrr}
\text{register}&H&k&u&A_\tau&A_g&Q&V&M&U&S,D\\\hline
\text{degree}&4&5&11&22&28&6&7&2&6&10
\end{array}
\]
in the ordinary case, with $\deg G=28$. The normalized case instead has
$\deg A_f=14$, $\deg Q=22$ and $\deg V=19$.
Direct degree propagation through the exact DAGs proves the final upper
bounds $378,172,98$ for ordinary fourteen, ordinary fifteen and centered
sixteen variables. Two cancellation-sensitive cases require more.

\paragraph{Thirteen variables.}
Pair the two signs of the auxiliary root first. Write
\[
 Z=UV+s,\qquad s=Un+Mb_0+H(v-q^2),
\]
and, for each of the sixteen sign choices, put
\[
 B=s\mathbin{\pm}9Uf\mathbin{\pm}3U\tau
     \mathbin{\pm}Md_0\mathbin{\pm}27Up_1.
\]
Assign formal root weights $7,11,5,14$ to $f,\tau,d_0,p_1$.
Each squared radicand has degree at most twice its root weight, so
reduction in this formal root algebra never increases weighted degree.
Thus $B$ has weight at most $20$. The paired factor is exactly
\begin{align*}
 \delta\{(UV+B)^2-U^2C_y/\delta\}
 &=\delta B^2+2\delta UVB-U^2(V^2-1).
\end{align*}
The identity uses $C_y=(\delta+1)V^2-1$. Its three terms have bounds
$62,67,50$. The sixteen-factor product therefore has degree at most
$16\cdot67=1072$, and is exactly the cleared five-root norm.
This argument is polynomial algebra, including all exceptional inputs.

\paragraph{Centered fifteen variables.}
Use the exact centered two-root identity with its scale $M$ replaced by
$U$, and its carrier argument replaced by
\[
 x=Un+2Mb_0+H(v-q^2)\mathbin{\pm}2Md_0.
\]
With $d_0$ of weight five, $x$ has weight at most seven. Define
\begin{align*}
 j&=\delta\{x^2+U(6t_0+4V)x+9U^2\ell\}
      -4U^2(V^2-1),\\
 h&=12\delta t_0VU^2.
\end{align*}
Their weights are at most $30$ and $36$. Every term of
\[
 j^2+2hj+576\delta U^4A_\tau(1-V^2)
       +144\delta^2U^4\ell V^2
\]
has weight at most $66$. The exact identity uses only
$t_0^2=4A_\tau+\ell$, not a Pell zero. Taking the two main-root
conjugates gives norm degree at most $132$; the strong and guard
finalizers add $10+28$, yielding $170$.

Independent modular-line certificates modulo $65521$ and $65519$, using
the actual six numeral byte strings, attain every stated bound with a
nonzero leading coefficient. They are rigorous lower certificates
combined with the preceding universal upper proofs.
All literal DAGs pass topology, variable closure, liveness and coefficient
hash checks. An independent exact sparse-algebra audit verifies every
prefix, carrier, radicand, norm and finalizer, using all $20\,349$
coefficients of the generic five-root norm; deliberate corruptions are
rejected. Signed modular diagnostics additionally include $c=0$, $v=0$
and their intersection. These algebraic checks do not establish an
unbounded compiler theorem: the named ordinary source calibration remains
an explicit mathematical dependency.
